\section{Interpretabilidad, depuración y automatización}

\subsection{Visualización y diagnóstico de los valores Q}

Durante el entrenamiento, se puede observar con un heatmap o un contour cómo se distribuyen los valores Q según cambian las acciones y los estados. Si una acción recibe siempre un puntaje Q muy alto y, aun así, nunca se elige, eso puede indicar que la implementación de la política de comportamiento (por ejemplo, $\epsilon$-greedy) está afectada por un bug.

\begin{knowledgebox}{Elementos de las herramientas de diagnóstico}
La información importante incluye: los gradientes de Q, la curva del TD error y la action frequency en el Replay buffer. Para mostrar gráficas de series de tiempo se pueden usar TensorBoard y MLFlow.
\end{knowledgebox}

\subsection{Automatización del pipeline de entrenamiento}

Se recomienda incorporar comprobaciones automáticas en el ciclo de entrenamiento: cada cierto número de epochs, hacer una offline evaluation, sincronizar las metrics con el dashboard y calcular automáticamente la sample efficiency. El pipeline de entrenamiento también debe admitir `checkpoint+resume` y `metric gating`; este último pausa el entrenamiento automáticamente cuando bajan las métricas clave.

\begin{importantbox}{Centinela del Replay buffer}
Un modelo ligero calcula en tiempo real la coverage y las reward statistics del buffer, y emite una alerta en cuanto detecta un desplazamiento brusco de la distribución. Este centinela permite evitar que, por recolectar datos one-hot desde las primeras etapas, el entrenamiento posterior colapse.
\end{importantbox}

\begin{warningbox}{Dependencia excesiva de la automatización en métricas históricas}
Si el gating solo considera el promedio histórico de la reward, puede pasar por alto el recent decay. Se recomienda introducir una sliding window version de las métricas y escalar las anomalías a un operator para que las confirme manualmente.
\end{warningbox}

\subsection{Resumen de la sección}

La interpretabilidad y la automatización, trabajando en paralelo, son la base del entrenamiento a gran escala; organizar bien las vistas de diagnóstico y los mecanismos automáticos de gating reduce de manera considerable el número de reinicios.
